\section{Circle palettes and their cardinality}\label{sec:palettes}

We now assume a weak measurable five-coloring and develop angular
palettes at every center.  Their labels record all possible limiting
circle samples.  The main result of this section is that almost every
palette contains at most two labels, at every fixed center.  We obtain
this bound from exclusions between circle samples and a separation
property for centers whose palettes meet one fixed pair of labels.

\subsection{Typical samples, traces, and maximal palettes}

For this section and the subsequent contradiction, the label set is
\(\{1,\ldots,5\}\); write \(A_i\) for the classes,
\(a_i=\ind_{A_i}\), and \(c\) for the coloring.
We choose Borel representatives forming a partition.
Changing a coloring on a plane-null set leaves the weak condition
unchanged, because that set is hit at either endpoint only on a
\(dx\,d\sigma(u)\)-null set.

Call a point \emph{typical} if it has density one of its actual color.
Typical points form a conull set.  By
Lemma~\ref{lem:palettes-density-proper}, any unit pair of typical
points has different colors.  The same lemma applies to two
density-one points of a class regardless of the values assigned
to the coloring at those individual points.
For each label let
\[
 H_i=\{q\in\R^2:
       \abs{A_i\cap B(q,r)}>0\text{ for every }r>0\}
\]
be its closed spatial essential support.

Let \(S=S^1\) carry uniform angular probability \(\sigma\).
Let \(R\) denote rotation by \(\kappa=\pi/3\), and \(J\) rotation
by \(\pi/2\).  We identify unit directions with angles modulo
\(2\pi\).  Two directions are called parallel if they are equal
or opposite; all angular density and essential-support statements
refer to angular Lebesgue measure.

Call \((y,r)\), with \(y\in\R^2\) and \(r>0\), a \emph{typical circle
sampling} if \(y+re\) is typical for almost every \(e\in S\).
For each fixed \(r\), this holds for almost every \(y\) by Fubini
and translation invariance of plane measure.
For each fixed \(y\), it holds for almost every \(r>0\) by polar
coordinates and Fubini.

\begin{definition}[Circle traces and maximal palettes]
\label{def:palettes-traces}
For each fixed center \(x\), let \(\mathcal W_x\) be the family of
vector functions in \(L^\infty(S;\R^5)\) that can be obtained as
\[
 b=\operatorname{w}^*\!-\!\lim_{n\to\infty}
          (a_i(y_n+r_ne))_{i=1}^5,\qquad
 y_n\to x,\quad r_n\to1,
\]
where every \((y_n,r_n)\) is a typical circle sampling.
Constant sequences are permitted.
For a label \(i\), form the essential measurable union
\[
 D_{x,i}=\mathop{\operatorname{ess\,union}}_{b\in\mathcal W_x}
                            \{e:b_i(e)>0\}.
\]
The maximal palette is the measurable set-valued function
\[
 P_x(e)=\{i:e\in D_{x,i}\},
\]
defined up to angular null sets.
\end{definition}

Here are the compactness and countability facts in this definition.
Typical samplings with radius one can be chosen at centers approaching
any prescribed \(x\).  Their bounded vectors have weak-* convergent
subsequences, since \(L^1(S)\) is separable and the bounded ball
of its dual is weak-* compact and metrizable.  Thus \(\mathcal W_x\)
is nonempty.  Every \(b\in\mathcal W_x\) has nonnegative coordinates
and \(\sum_i b_i=1\) almost everywhere, since these assertions pass
to a weak-* limit by integration against nonnegative test functions.

The family also has the following closure under varying centers:
\begin{equation}\label{eq:palettes-trace-closure}
 x_n\to x,\quad b^n\in\mathcal W_{x_n},\quad
 b^n\rightharpoonup^* b
 \quad\Longrightarrow\quad b\in\mathcal W_x .
\end{equation}
To see this, fix a metric for the bounded-ball weak-* topology.
From a sequence realizing \(b^n\), choose one sample within \(1/n\)
of \(b^n\) in that metric, with center within \(1/n\) of \(x_n\)
and radius within \(1/n\) of one.  These chosen samples realize \(b\)
at \(x\).  In particular \(\mathcal W_x\) is compact.

For completeness, the essential union for a fixed label can always
be taken over a countable subfamily.  Let \(M_i\) be the supremum
of the measures of finite unions of the sets \(\{b_i>0\}\).
Choose finite unions approaching \(M_i\), and let \(D_{x,i}\) be
their countable union.  Its measure is \(M_i\).  If another member
\(\{b_i>0\}\) had positive measure outside this union, adjoining it
to sufficiently large finite partial unions would contradict the
definition of \(M_i\).  This proves that every individual trace
has its positive coordinates in the palette almost everywhere.
It also gives nonemptiness of \(P_x(e)\) almost everywhere, by
using any one trace and its sum-one property.

The null set in the last assertion may depend on the individual
trace, and the angular null sets may depend on \(x\).  We shall use
countable generating trace families whenever a conclusion is to be
asserted for the whole palette.  No joint measurability of \(P_x\)
in its center is assumed.

\subsection{Exclusion, density, and smooth transport}

\begin{lemma}[Palette transfer]\label{lem:palette-transfer}
The following assertions hold at every fixed center \(x\).

\begin{enumerate}
\item For almost every independent-angle pair \((e,v)\in S\times S\),
\begin{equation}\label{eq:source-10}
 c(x+e+v)\notin P_x(e)\cup P_x(v).
\end{equation}

\item If \(e\) is an angular density-one point of
\(D_{x,i}=\{e':i\in P_x(e')\}\), then for each \(v\) nonparallel
to \(e\), the spatial density of \(A_i\) at \(x+e+v\) is zero.

\item Let \(z\) be any fixed center, possibly \(z=x\).  Suppose on an open angular
arc, or locally on a union of such arcs, there is a smooth representation
\[
 z+f=x+s(f)+t(f),\qquad s(f),t(f)\in S,
\]
with \(s(f),t(f)\) nonparallel and \(s\) a local diffeomorphism in
the angular coordinate.  Then almost everywhere on that domain,
\begin{equation}\label{eq:source-11}
 P_z(f)\cap P_x(s(f))=\varnothing.
\end{equation}
\end{enumerate}
Each assertion permits an exceptional angular set depending on the
fixed centers and on the specified maps.
\end{lemma}

\begin{proof}
Fix \(b\in\mathcal W_x\), with realizing samplings \((y_n,r_n)\).
For every \(n\), the point \(y_n+r_ne\) is typical for almost every
\(e\).  The second point \(y_n+r_ne+v\) is typical for almost every
pair \((e,v)\): in angular coordinates the map
\[
 G_n(e,v)=y_n+r_ne+v
\]
has Jacobian of absolute value
\(r_n\abs{\sin(\arg v-\arg e)}\), which is nonsingular off a
product-null set.  On countably many smaller coordinate patches
away from that set, change of variables pulls back any plane-null
set to a product-null set.  The two typical points have distance
one.  Lemma~\ref{lem:palettes-density-proper} therefore gives
\begin{equation}\label{eq:palettes-sample-disjoint}
 a_i(y_n+r_ne)a_i(y_n+r_ne+v)=0
 \quad\text{for almost every }(e,v).
\end{equation}

We need the following strong convergence for the second factor:
\begin{equation}\label{eq:palettes-strong-circle-sum}
 a_i(y_n+r_ne+v)\longrightarrow a_i(x+e+v)
 \quad\text{in }L^1(S\times S).
\end{equation}
Here is a change-of-variables proof which also covers arbitrary
measurable color boundaries.  For \(\delta>0\), discard the angular
pairs with
\(\abs{\sin(\arg v-\arg e)}\le\delta\).  Their measure tends to
zero with \(\delta\).  On the remaining set, for large \(n\), the
Jacobian is bounded below by \(\delta/2\).  The multiplicity of
\(G_n\) there is at most two: an inverse image of a specified
sum amounts to an intersection of the circle of radius \(r_n\)
with a translated circle of radius one.  A coincident-circle
case can occur only at the antiparallel, hence discarded, directions.
The same assertions hold for \(G(e,v)=x+e+v\).
Their images lie in one fixed compact plane region for all large \(n\).

It follows that on the retained angular set the pushforwards of
angular product measure have densities bounded by a constant
depending only on \(\delta\), uniformly in \(n\).
Approximate \(a_i\) on a compact neighborhood of those images by
a continuous function taking values in \([0,1]\), in plane \(L^1\).
The density bounds control its composition errors uniformly.
For the continuous approximant the compositions converge uniformly,
since \(y_n\to x\) and \(r_n\to1\).
First making the approximation error small, then letting \(n\to\infty\),
and finally letting \(\delta\downarrow0\), proves
\eqref{eq:palettes-strong-circle-sum}.

Set \(B_n(e)=a_i(y_n+r_ne)\) and
\(Q(e,v)=a_i(x+e+v)\).  The error in replacing the second factor
of \eqref{eq:palettes-sample-disjoint} by \(Q\), after integration,
is at most the \(L^1\) error in
\eqref{eq:palettes-strong-circle-sum}.  Weak-* convergence of \(B_n\)
then applies to the fixed \(L^1(S)\) test
\(\int_S Q(e,v)\,d\sigma(v)\).  We obtain
\[
 \int_{S\times S}b_i(e)a_i(x+e+v)\,d\sigma(e)\,d\sigma(v)=0.
\]
Both factors are nonnegative, so their product is zero almost
everywhere.  Taking the countable generating trace families in
Definition~\ref{def:palettes-traces}, for all five labels, proves
exclusion of \(P_x(e)\).  Interchanging \(e,v\) proves exclusion
of \(P_x(v)\), and gives \eqref{eq:source-10}.
In particular this passage used a strongly convergent factor
against a weak-* limit, not a product of two weak-* limits.

Next fix an angular density-one point \(e\) of \(D_{x,i}\) and
any nonparallel direction \(v\).  The map
\[
 (e',v')\longmapsto x+e'+v'
\]
is a smooth local diffeomorphism near \((e,v)\).
Choose real angular coordinates near \(e,v\) and an inverse branch
on a plane neighborhood of \(x+e+v\).
For small \(h\), the inverse image under this branch of
\(B(x+e+v,h)\) lies in a product of angular intervals of length
\(O(h)\) about \(e,v\), and its change-of-variables factors are
bounded.  On this branch, \eqref{eq:source-10} implies that a
point of \(A_i\) has first coordinate outside \(D_{x,i}\), apart
from a spatial null set.  Therefore
\[
 \abs{A_i\cap B(x+e+v,h)}
 \le C h\,
       \abs{[e-Ch,e+Ch]\setminus D_{x,i}}
 =o(h^2).
\]
The constants can depend on the fixed nonparallel pair.  This
establishes assertion 2 for each such \(v\).  It is a
two-dimensional neighborhood argument and does not restrict
\eqref{eq:source-10} to a possibly exceptional single-angle fiber.

Finally fix the centers and representation in assertion 3, and work
first on a smaller compact arc on which one inverse branch of the
two-angle sum map is available.  Let a trace \(b\in\mathcal W_z\)
be realized by \((z_n,r_n)\).  Applying this inverse branch to
\(z_n+r_nf-x\) gives perturbed representations
\[
 z_n+r_nf=x+s_n(f)+t_n(f),
\]
where \(s_n,t_n\) are nonparallel and converge to \(s,t\) in \(C^1\)
on the smaller arc.  Shrinking the arc if necessary, \(s\) is
a diffeomorphism with derivative bounded away from zero, and the
\(s_n\) have uniformly bounded inverse derivatives for large \(n\).

For \(D_i=D_{x,i}\), these facts imply
\begin{equation}\label{eq:palettes-pullback-continuity}
 \ind_{D_i}(s_n(f))\longrightarrow\ind_{D_i}(s(f))
 \quad\text{in }L^1\text{ on the smaller arc}.
\end{equation}
To verify it, approximate \(\ind_{D_i}\) on a common image interval
by a continuous function in angular \(L^1\).  Change of variables
and the inverse-derivative bounds control the errors of composition
uniformly for \(s_n\) and for \(s\).  Uniform convergence of \(s_n\)
handles the continuous approximant.

Almost every \(f\) with \(s_n(f)\in D_i\) has \(s_n(f)\) a
density-one point of \(D_i\).  Indeed the set of exceptions in
\(D_i\) is angular null by differentiation, and \(s_n\) pulls
it back to a null set.  Assertion 2 makes the density of \(A_i\)
at \(z_n+r_nf\) zero for all these \(f\).  The sampling
\((z_n,r_n)\) is typical on its circle, so for almost every such
\(f\) its actual color cannot be \(i\).  Consequently
\[
 a_i(z_n+r_nf)\ind_{D_i}(s_n(f))=0
 \quad\text{almost everywhere on the smaller arc}.
\]
Use \eqref{eq:palettes-pullback-continuity} to replace the second
factor strongly in \(L^1\), and then use weak-* convergence of the
first factor.  Nonnegativity gives
\[
 b_i(f)\ind_{D_i}(s(f))=0
 \quad\text{almost everywhere}.
\]
Taking a countable generating family of traces at \(z\) and the
five labels proves \eqref{eq:source-11} on this arc.
A countable cover by smaller arcs proves it on the stated domain.
\end{proof}

\begin{remark}\label{rem:palettes-transport-scope}
In assertion 3 of Lemma~\ref{lem:palette-transfer}, the second endpoint
map \(t\) may be constant.  Nonparallelness makes the two-dimensional
sum map invertible, while the local diffeomorphism requirement on
\(s\) alone controls preimages of angular null sets.
In particular, for any fixed direction \(a\), the representation
\[
 (x+a)+f=x+f+a
\]
is admissible off \(f=\pm a\).  The proof applies at every fixed
pair of centers and never needs a common exceptional angular set
over an uncountable family of centers.
\end{remark}

\begin{corollary}[Sixty-degree exclusion]
\label{cor:palettes-sixty-degree}
At every fixed center \(x\),
\begin{equation}\label{eq:source-12}
 P_x(f)\cap P_x(Rf)=\varnothing
 \qquad\text{for almost every }f\in S.
\end{equation}
\end{corollary}

\begin{proof}
Use \(x+f=x+Rf+R^{-1}f\) in
Lemma~\ref{lem:palette-transfer}.  The endpoint directions are
nonparallel and the first endpoint map is a rotation.
\end{proof}

\subsection{Separation of binary centers}
\label{subsec:palettes-binary}

Fix two distinct labels \(i_+,i_-\).  Call a center \(x\)
\emph{binary for \(\{i_+,i_-\}\)} if
\[
 P_x(e)\cap\{i_+,i_-\}\ne\varnothing
 \qquad\text{for almost every }e\in S.
\]
The exclusion \eqref{eq:source-12} makes this intersection a singleton
almost everywhere.  Indeed, at almost every \(e\), the intersections
at \(e\) and \(Re\) are nonempty disjoint subsets of a two-element set;
they must therefore be the two different singletons.  Thus there is a
unique measurable selection from this pair, up to a null set, and it
alternates under \(R\).  Conversely, the existence of an alternating
selection from the pair plainly implies the displayed intersection
condition.

Write this selection as a measurable \(2\pi\)-periodic sign function
\(s_x:\R\to\{-1,1\}\), assigning the signs \(+1,-1\) to \(i_+,i_-\)
respectively and choosing arbitrary values on the exceptional null
set.  With \(\kappa=\pi/3\), it satisfies
\begin{equation}\label{eq:palettes-binary-signs}
 s_x(t+\kappa)=-s_x(t)
 \qquad\text{for almost every }t.
\end{equation}
In particular, either sign occupies angular measure \(\pi\) in a
period.

\begin{lemma}[Binary-center separation]\label{lem:binary-separation}
There is an absolute constant \(\delta_*>0\) such that, for every
fixed pair of labels, any two distinct centers binary for that pair
are at distance at least \(\delta_*\).
\end{lemma}

\begin{proof}
Suppose that \(x,y\) are distinct binary centers for the same pair,
and put \(d=|y-x|\).  We will derive a contradiction whenever
\(0<d<\delta_*\), with \(\delta_*\) chosen below independently of the
coloring, the pair, and the centers.  Rotate angular coordinates so
that \(y-x=d\) lies on the positive real axis.

\smallskip
\noindent\emph{The maps forced by palette transport.}
For \(d<1/2\), the vector \(d+e^{it}\) has length strictly between
zero and two.  Define real lifts
\[
 \begin{split}
 \Theta_d(t)&=t+\operatorname{Im}\log(1+d e^{-it}),\\
 \alpha_\pm(t)&=\Theta_d(t)
       \pm\arccos\bigl(|d+e^{it}|/2\bigr),
 \end{split}
\]
where the logarithm is the branch near \(1\).  These formulas give
the exact representation
\[
 y+e^{it}
   =x+e^{i\alpha_+(t)}+e^{i\alpha_-(t)}.
\]
Its two unit directions are nonparallel.  Their angular difference
is twice an angle strictly between zero and \(\pi/2\), and stays
away from both parallel cases for sufficiently small \(d\).

The functions \(\alpha_\pm\) are smooth degree-one lifts:
\(\alpha_\pm(t+2\pi)=\alpha_\pm(t)+2\pi\).  Their derivatives converge
uniformly to \(1\) as \(d\to0\), so for small \(d\) they induce
orientation-preserving circle diffeomorphisms.  We may therefore
apply Lemma~\ref{lem:palette-transfer} locally, with either endpoint
as its first direction, and cover the angular circle by finitely
many such domains.  It follows that
\[
 P_y(e^{it})\cap P_x(e^{i\alpha_\pm(t)})=\varnothing
 \qquad\text{for almost every }t.
\]
Both binary selections belong to the same fixed pair of labels.
Consequently,
\[
 s_x(\alpha_+(t))=s_x(\alpha_-(t))=-s_y(t)
 \qquad\text{for almost every }t.
\]
Since \(\alpha_-\) is nonsingular for angular Lebesgue measure,
\(s_x\) is invariant almost everywhere under
\(H=\alpha_+\circ\alpha_-^{-1}\).

By \eqref{eq:palettes-binary-signs}, \(s_x\) is also invariant under
translation by \(2\kappa\).  Define
\[
 F(t)=H(t)-2\kappa,\qquad
 F_j(t)=j\kappa+F(t-j\kappa),\qquad 0\le j\le5.
\]
Then \(s_x\circ F=s_x\) almost everywhere.  More generally,
\eqref{eq:palettes-binary-signs} gives
\(s_x(t+j\kappa)=(-1)^j s_x(t)\) almost everywhere, and hence
\[
 s_x(F_j(t))
 =(-1)^j s_x(F(t-j\kappa))
 =(-1)^{2j}s_x(t)
 =s_x(t)
\]
almost everywhere.  The changes of variables in these identities
are legitimate because all the maps involved are diffeomorphisms.

\smallskip
\noindent\emph{Uniform expansions of the six maps.}
For a periodic function, an \(O_{C^2}(d^2)\) remainder means that
its value and its first two derivatives in \(t\) are bounded by
a fixed constant times \(d^2\), uniformly in \(t\).  Taylor expansion
gives
\[
 \begin{split}
 \operatorname{Im}\log(1+d e^{-it})
     &=-d\sin t+O_{C^2}(d^2),\\
 |d+e^{it}|&=1+d\cos t+O_{C^2}(d^2),\\
 \alpha_\pm(t)
     &=t\pm\kappa-d\sin t
            \mp\frac d{\sqrt3}\cos t+O_{C^2}(d^2).
 \end{split}
\]
All the bounds are uniform: the logarithm, square root, and
arccosine in these formulas range in fixed smooth domains for
\(d<1/2\).  The inverse maps also depend smoothly on \(d\), with
uniformly bounded derivatives on a period.

For the composition, put
\[
 a_\pm(t)=-\sin t\mp\frac1{\sqrt3}\cos t.
\]
Then
\[
 \alpha_-^{-1}(t)
   =t+\kappa-d a_-(t+\kappa)+O_{C^2}(d^2),
\]
and therefore
\begin{equation}\label{eq:palettes-binary-map-expansion}
 \begin{split}
 F(t)
 &=t+d\bigl(a_+(t+\kappa)-a_-(t+\kappa)\bigr)
                                      +O_{C^2}(d^2)\\
 &=t-\frac{2d}{\sqrt3}\cos(t+\kappa)+O_{C^2}(d^2)\\
 &=t+\frac{2d}{\sqrt3}\sin(t-\pi/6)+O_{C^2}(d^2).
 \end{split}
\end{equation}
Thus
\[
 F_j(t)=t+\frac{2d}{\sqrt3}
             \sin(t-j\kappa-\pi/6)+O_{C^2}(d^2).
\]
Among six phases spaced by \(\pi/3\), the maximum sine and the
maximum cosine are each at least \(\sqrt3/2\).  Absorbing the
quadratic remainders, we obtain absolute constants
\(c_0,C_0,d_1>0\), with \(d_1<1/2\), such that whenever
\(0<d<d_1\), for every \(t\),
\begin{equation}\label{eq:palettes-binary-control}
 \begin{split}
 \max_{0\le j\le5}(F_j(t)-t)&\ge c_0d,\\
 \max_{0\le j\le5}F_j'(t)&\ge1+c_0d,\\
 |F_j(t)-t|,\quad |F_j'(t)-1|,\quad |F_j''(t)|
     &\le C_0d\qquad(0\le j\le5).
 \end{split}
\end{equation}
The map attaining the first maximum need not attain the second.
We will use the derivative maximum to enlarge small intervals,
and the displacement maximum later to extend a conull interval.

Choose once and for all
\[
 0<L\le\min\{c_0/C_0,\pi/4\},
 \qquad
 \delta_*=\min\{d_1,(2C_0)^{-1},L/(2C_0)\},
\]
and assume \(0<d<\delta_*\).  In particular every \(F_j\) and its
inverse is Lipschitz, since
\[
 \frac12\le F_j'\le\frac32.
\]

\smallskip
\noindent\emph{Enlarging density intervals with bounded distortion.}
This is an adaptation of Sullivan's expanding-action and bounded-distortion
strategy; see Deroin--Kleptsyn--Navas \cite[Section~2.1]{DeroinKleptsynNavas2008}.
The argument below replaces the usual minimality step by explicit
small-displacement overlap propagation when extending the conull interval.
Let
\[
 A=\{t\in\R:s_x(t)=1\}.
\]
It is measurable and \(2\pi\)-periodic, and has measure \(\pi\)
in every period.  The identities already proved say that
\(F_j^{-1}(A)\) and \(A\) differ by a null set.  Because \(F_j\)
and its inverse preserve null sets, the same holds for
\(F_j(A)\) and \(A\), and for every finite composition of the maps.
There is no difficulty from choosing the compositions adaptively:
there are only countably many finite words in the six maps and
their inverses.  Equivalently, one may discard the union of the
images of the finitely many exceptional null sets under this
countable group, obtaining a common invariant conull set on which
all the membership identities hold.

Take an interval \(I_0\) of positive length \(\ell_0<L\).
Given \(I_r\) with length \(\ell_r<L\), choose \(F_{j_r}\) whose
derivative at the midpoint of \(I_r\) is at least \(1+c_0d\), and
put \(I_{r+1}=F_{j_r}(I_r)\).  For every \(t\in I_r\),
\[
 F_{j_r}'(t)
 \ge1+c_0d-\frac12C_0d\ell_r
 \ge1+\frac12c_0d.
\]
Writing \(\lambda=1+c_0d/2>1\), we therefore have
\(\ell_{r+1}\ge\lambda\ell_r\).  The process stops after finitely
many steps, at the first \(N\) for which \(\ell_N\ge L\), and
\[
 L\le\ell_N\le(1+C_0d)L.
\]

Let \(G=F_{j_{N-1}}\circ\cdots\circ F_{j_0}\).
For each elementary map,
\[
 |(\log F_j')'|=\left|\frac{F_j''}{F_j'}\right|\le2C_0d.
\]
Applying the chain rule at two points of \(I_0\), whose images
before step \(r\) lie in \(I_r\), gives
\[
 \operatorname{osc}_{I_0}\log G'
 \le2C_0d\sum_{r=0}^{N-1}\ell_r.
\]
The geometric growth bounds the sum independently of \(N\):
\[
 \sum_{r=0}^{N-1}\ell_r
 \le\ell_N\sum_{q=1}^{N}\lambda^{-q}
 \le\frac{\ell_N}{\lambda-1}.
\]
Consequently,
\begin{equation}\label{eq:palettes-binary-distortion}
 \operatorname{osc}_{I_0}\log G'
 \le\frac{4C_0(1+C_0d)L}{c_0}
 \le\frac{6C_0L}{c_0}.
\end{equation}
In particular,
\[
 \frac{\sup_{I_0}G'}{\inf_{I_0}G'}
 \le D_{\rm dist},\qquad D_{\rm dist}=\exp(6C_0L/c_0),
\]
uniformly in the initial interval, the number of steps, and
\(0<d<\delta_*\).

The invariance of \(A\), together with preservation of null sets,
implies that \(G(I_0)\setminus A\) and \(G(I_0\setminus A)\)
differ by a null set.  Change of variables and the distortion
bound therefore give
\[
 \frac{|I_N\setminus A|}{|I_N|}
 \le
 \frac{\sup_{I_0}G'}{\inf_{I_0}G'}
       \frac{|I_0\setminus A|}{|I_0|}
 \le D_{\rm dist}\,\frac{|I_0\setminus A|}{|I_0|}.
\]
This is the only measure comparison needed; the maps need not
preserve Lebesgue measure.

Since \(A\) has positive measure, choose a Lebesgue density point
of \(A\) and a sequence of centered intervals shrinking to it.
Their relative complements in \(A\) tend to zero.  Apply the
construction to each interval.  The resulting intervals \(J_n\)
have lengths in \([L,(1+C_0d)L]\), and
\[
 |J_n\setminus A|\longrightarrow0.
\]
Translate \(J_n\) by an integer multiple of \(2\pi\) so that its
left endpoint lies in \([0,2\pi)\).  This leaves the preceding
measure statement unchanged.  Passing to a subsequence, their
two endpoints converge to those of an interval \(J=[a,b]\) of
length at least \(L\).  Endpoint convergence gives
\(|J_n\mathbin{\triangle}J|\to0\), and hence
\[
 |J\setminus A|
 \le |J\mathbin{\triangle}J_n|+|J_n\setminus A|
 \longrightarrow0.
\]
We have obtained an interval of fixed positive length on which
\(A\) is conull.

\smallskip
\noindent\emph{Extending a conull interval.}
Suppose \(A\) is conull on \([a,b]\), with \(b-a\ge L\).
By the first inequality in \eqref{eq:palettes-binary-control},
choose \(j\) so that
\[
 F_j(b)\ge b+c_0d.
\]
The image \(F_j([a,b])\) is also conull in \(A\).
Moreover,
\[
 F_j(a)\le a+C_0d<a+L\le b.
\]
Thus this image overlaps \([a,b]\), and their union contains
\([a,F_j(b)]\), on which \(A\) is conull.  Its right endpoint
has advanced by at least \(c_0d\).  This argument concerns
geometric endpoints only; no membership in \(A\) is asserted
at either endpoint.

Repeat, retaining the left endpoint \(a\).  After finitely
many steps the conull interval has length at least \(2\pi\).
Periodicity then forces \(A\) to have full measure in a period,
contrary to the measure \(\pi\) forced by
\eqref{eq:palettes-binary-signs}.  This contradiction excludes
\(0<|y-x|<\delta_*\) and proves the lemma.
\end{proof}

\subsection{At most two labels in an angular palette}

\begin{proposition}[Palette cardinality]\label{prop:palette-size}
For every center \(x\in\R^2\),
\begin{equation}\label{eq:source-13}
 1\le\abs{P_x(e)}\le2
 \qquad\text{for almost every }e\in S.
\end{equation}
The angular exceptional set may depend on \(x\).
\end{proposition}

\begin{proof}
Nonemptiness follows from any one trace in \(\mathcal W_x\), whose
coordinates are nonnegative and sum to one.
Suppose the upper bound fails on a positive-measure set of angles.
There are only finitely many triples of labels, so some fixed triple
\(\mathcal A\subset\{1,\ldots,5\}\) is contained in \(P_x(e)\) on a
positive-measure angular set \(B\).

By \eqref{eq:source-10} and Fubini, for almost every \(e\in B\),
the actual coloring on the circle
\[
 v\longmapsto x+e+v
\]
uses only the complementary pair
\(\{1,\ldots,5\}\setminus\mathcal A\) almost everywhere.
For almost every such \(e\), this exact radius-one circle sampling
is also typical.  Indeed the map
\((e,v)\mapsto x+e+v\) is nonsingular off a product-null set,
so the inverse image of the plane-null set of atypical points
is product-null; Fubini then gives the asserted slices.

For each of these \(e\), the actual sampling at center \(x+e\),
radius exactly one, is itself a trace in \(\mathcal W_{x+e}\),
using a constant sequence.  Its coordinates sum to one and are
supported in the complementary pair.  Hence the maximal palette
at \(x+e\) intersects that pair almost everywhere.  By
\eqref{eq:source-12}, it has a unique member of the pair almost
everywhere, alternating at an \(R\)-step.  Thus \(x+e\) is binary
for that same fixed pair, even if other traces add further labels
to the maximal palette.

The good angles still have positive measure, so they give infinitely
many distinct binary centers on the compact circle \(x+S\).
Lemma~\ref{lem:binary-separation} makes these centers uniformly
separated.  A uniformly separated subset of a compact metric space
is finite, by a finite cover with sufficiently small balls.
This contradiction proves the upper bound.
\end{proof}
